\documentclass{beamer}
\usepackage{listings,xcolor,epsf,mathrsfs}
\usepackage{graphicx}
\usepackage{booktabs}
\usepackage{array}
\usepackage{amsmath}
\usepackage{multirow}
\usetheme{Madrid}

\newcommand{\pa}{\partial}
\newcommand{\di}{\displaystyle}
\newcommand{\acc}{\mathrm{Accuracy}}

\begin{document}
 \mode<presentation>

\title[BO-SVM EEG Classification]{EEG-Based Epileptic Signal Classification Using BO-SVM and Hierarchical SVM}
\author{Your Name}
\date{\today}
\institute[XMUM]{Department of Mathematics\\ Xiamen University Malaysia}

\begin{frame}
\titlepage
\end{frame}

\begin{frame}
\frametitle{Presentation Outline}
\begin{enumerate}
\item Background and dataset
\item Preprocessing and feature extraction
\item BO-SVM baseline
\item Error analysis
\item Hierarchical SVM improvement
\item Experimental results
\item Discussion and conclusion
\end{enumerate}
\end{frame}

\begin{frame}
\frametitle{Background}
\framesubtitle{\it{Automatic EEG classification for epilepsy analysis}}
\begin{itemize}
\item Epilepsy is a neurological disorder characterized by abnormal brain activity.
\item Electroencephalography (EEG) records brain electrical signals directly.
\item Automatic EEG classification can reduce manual workload and support clinical decision-making.
\item This project develops a feature-based SVM pipeline for Bonn EEG classification.
\end{itemize}
\end{frame}

\begin{frame}
\frametitle{Bonn EEG Dataset}
\begin{itemize}
\item Dataset: University of Bonn EEG dataset.
\item Five classes, with 100 EEG segments per class.
\item Sampling frequency: 173.61 Hz.
\item Each EEG segment contains 4097 samples.
\end{itemize}

\begin{table}
\centering
\begin{tabular}{cl}
\toprule
Class & Description\\
\midrule
Z & Healthy subject, eyes open\\
O & Healthy subject, eyes closed\\
N & Interictal EEG from epileptogenic zone\\
F & Interictal EEG from non-epileptogenic zone\\
S & Ictal / seizure EEG\\
\bottomrule
\end{tabular}
\end{table}
\end{frame}

\begin{frame}
\frametitle{Project Objectives}
\begin{enumerate}
\item Build an EEG classification pipeline for the Bonn dataset.
\item Compare SVM without optimization and Bayesian-Optimized SVM.
\item Identify the main source of classification errors.
\item Improve performance using hierarchical SVM.
\item Evaluate both five-class and subset classification tasks.
\end{enumerate}
\end{frame}

\begin{frame}
\frametitle{Overall Methodology}
\framesubtitle{\it{Feature-based machine learning pipeline}}
\begin{center}
\begin{tabular}{c}
Raw EEG signal\\
$\Downarrow$\\
Band-pass filtering\\
$\Downarrow$\\
Window segmentation\\
$\Downarrow$\\
Feature extraction\\
$\Downarrow$\\
Feature scaling\\
$\Downarrow$\\
RBF-SVM classification\\
$\Downarrow$\\
Evaluation
\end{tabular}
\end{center}
\end{frame}

\begin{frame}
\frametitle{Signal Preprocessing}
\begin{itemize}
\item Band-pass filtering was applied to retain useful EEG frequency components.
\item EEG segments were divided into fixed-length windows.
\item Non-overlapping windows were used for conservative evaluation.
\item 50\% overlapping windows were additionally tested to improve window-level performance.
\end{itemize}

\vspace{0.2cm}
\begin{block}{Evaluation Note}
The 89.20\% result uses non-overlapping windows, while the 90.33\% result uses 50\% overlapping windows.
\end{block}
\end{frame}

\begin{frame}
\frametitle{Feature Extraction}
\begin{itemize}
\item Frequency-domain features:
\begin{itemize}
\item PSD band power
\item Relative band power
\item Band power ratios
\end{itemize}
\item Statistical features:
\begin{itemize}
\item Mean, standard deviation, variance, RMS, line length
\item Skewness and kurtosis
\end{itemize}
\item Nonlinear features:
\begin{itemize}
\item Shannon entropy, spectral entropy, approximate entropy
\item Hjorth activity, mobility, and complexity
\end{itemize}
\item DWT features:
\begin{itemize}
\item Wavelet energy, relative wavelet energy, and wavelet entropy
\end{itemize}
\end{itemize}
\end{frame}

\begin{frame}
\frametitle{PSD Normalization}
\begin{itemize}
\item PSD describes how signal power is distributed across frequency bands.
\item Absolute band power may be affected by the overall signal amplitude.
\item Therefore, relative band power was computed:
\end{itemize}

$$
\text{Relative Power}_i =
\frac{\text{Band Power}_i}
{\sum_{j=1}^{5}\text{Band Power}_j}.
$$

\begin{block}{Purpose}
PSD normalization helps the model focus on the frequency distribution pattern rather than only the absolute signal magnitude.
\end{block}
\end{frame}

\begin{frame}
\frametitle{Discrete Wavelet Transform}
\begin{itemize}
\item EEG signals are non-stationary, so time-frequency features are useful.
\item DWT decomposes each EEG window into multiple wavelet sub-bands.
\item In this project, db4 wavelet decomposition was used.
\item From wavelet coefficients, energy, relative energy, entropy, mean, and standard deviation were extracted.
\end{itemize}

\begin{block}{Reason for Using DWT}
Enhanced DWT features were mainly used for N/F refinement, because N and F were the most difficult classes to separate.
\end{block}
\end{frame}

\begin{frame}
\frametitle{BO-SVM Baseline}
\begin{itemize}
\item RBF-SVM was selected as the base classifier.
\item SVM is suitable for small datasets and handcrafted high-dimensional features.
\item Bayesian Optimization was used to tune important hyperparameters.
\end{itemize}

\begin{table}
\centering
\begin{tabular}{ll}
\toprule
Parameter & Meaning\\
\midrule
$f_{low}$, $f_{high}$ & Band-pass frequency range\\
Window size & EEG window length\\
$C$ & SVM regularization parameter\\
$\gamma$ & RBF kernel parameter\\
Feature setting & Feature combination\\
\bottomrule
\end{tabular}
\end{table}
\end{frame}

\begin{frame}
\frametitle{Error Analysis}
\begin{itemize}
\item Initial BO-SVM achieved around 84--86\% test accuracy.
\item Confusion matrix analysis showed that class S was detected very well.
\item Most remaining errors came from confusion between N and F.
\end{itemize}

\begin{block}{Key Observation}
The main limitation was not seizure detection, but separating similar interictal EEG classes N and F.
\end{block}
\end{frame}

\begin{frame}
\frametitle{Improved Method: Hierarchical SVM}
\begin{itemize}
\item A single five-class SVM must learn all class boundaries at once.
\item To reduce N/F confusion, the task was decomposed into simpler stages.
\end{itemize}

\vspace{0.2cm}
\begin{center}
\begin{tabular}{cl}
\toprule
Stage & Classification task\\
\midrule
1 & S vs non-S\\
2 & Z/O/N/F classification\\
3 & N vs F refinement\\
\bottomrule
\end{tabular}
\end{center}

\begin{block}{Motivation}
A dedicated N/F classifier was added to handle the most difficult pair.
\end{block}
\end{frame}

\begin{frame}
\frametitle{Final Model Setting}
\begin{table}
\centering
\begin{tabular}{ll}
\toprule
Component & Setting\\
\midrule
Classifier & RBF-SVM\\
Stage design & S vs non-S, Z/O/N/F, N vs F\\
Main feature setting & entropy\_hjorth\\
N/F feature setting & nf\_enhanced\\
Feature selection & SelectKBest with mutual information\\
Hyperparameter search & GridSearchCV for each stage\\
Conservative evaluation & Non-overlapping windows\\
Enhanced evaluation & 50\% overlapping windows\\
\bottomrule
\end{tabular}
\end{table}
\end{frame}

\begin{frame}
\frametitle{Five-Class Classification Results}
\begin{table}
\centering
\begin{tabular}{lc}
\toprule
Method & Test Accuracy\\
\midrule
BO-SVM baseline & around 84--86\%\\
Hierarchical SVM, no overlap & \textbf{89.20\%}\\
Hierarchical SVM, 50\% overlap & \textbf{90.33\%}\\
\bottomrule
\end{tabular}
\end{table}

\begin{block}{Interpretation}
The hierarchical SVM improved performance by directly targeting the N/F confusion problem.
\end{block}
\end{frame}

\begin{frame}
\frametitle{Best Five-Class Result}
\framesubtitle{\it{Hierarchical SVM with 50\% overlapping windows}}
\begin{table}
\centering
\begin{tabular}{lc}
\toprule
Metric & Value\\
\midrule
Test accuracy & \textbf{90.33\%}\\
Macro recall & \textbf{90.33\%}\\
Recall Z & 93.33\%\\
Recall O & 89.44\%\\
Recall N & 83.33\%\\
Recall F & 86.67\%\\
Recall S & 98.89\%\\
\bottomrule
\end{tabular}
\end{table}
\end{frame}

\begin{frame}
\frametitle{Confusion Matrix: Five-Class Best Result}
\begin{center}
\begin{tabular}{c|ccccc}
\toprule
True / Pred & Z & O & N & F & S\\
\midrule
Z & 168 & 10 & 1 & 1 & 0\\
O & 18 & 161 & 1 & 0 & 0\\
N & 2 & 0 & 150 & 28 & 0\\
F & 4 & 0 & 20 & 156 & 0\\
S & 0 & 1 & 0 & 1 & 178\\
\bottomrule
\end{tabular}
\end{center}

\begin{block}{Observation}
S was classified very accurately. Most remaining errors were still between N and F.
\end{block}
\end{frame}

\begin{frame}
\frametitle{Subset Classification Experiments}
\framesubtitle{\it{Non-overlapping windows}}
\begin{table}
\centering
\begin{tabular}{llc}
\toprule
Task & Best Feature & Test Accuracy\\
\midrule
Z vs S & entropy\_hjorth & \textbf{100.00\%}\\
F vs S & nf\_enhanced & \textbf{99.00\%}\\
Z vs F vs S & nf\_enhanced & \textbf{98.67\%}\\
\bottomrule
\end{tabular}
\end{table}

\begin{block}{Conclusion}
Subset tasks show that S is easy to distinguish from Z and F. The full five-class task is harder due to similar non-seizure classes.
\end{block}
\end{frame}

\begin{frame}
\frametitle{Subset Confusion Matrices}
\begin{columns}
\begin{column}{0.31\textwidth}
\centering
\textbf{Z vs S}

\vspace{0.1cm}
\begin{tabular}{c|cc}
 & Z & S\\
\hline
Z & 100 & 0\\
S & 0 & 100\\
\end{tabular}
\end{column}

\begin{column}{0.31\textwidth}
\centering
\textbf{F vs S}

\vspace{0.1cm}
\begin{tabular}{c|cc}
 & F & S\\
\hline
F & 100 & 0\\
S & 2 & 98\\
\end{tabular}
\end{column}

\begin{column}{0.34\textwidth}
\centering
\textbf{Z vs F vs S}

\vspace{0.1cm}
\begin{tabular}{c|ccc}
 & Z & F & S\\
\hline
Z & 97 & 3 & 0\\
F & 0 & 99 & 1\\
S & 0 & 0 & 100\\
\end{tabular}
\end{column}
\end{columns}
\end{frame}

\begin{frame}
\frametitle{Comparison With Existing Studies}
\begin{itemize}
\item Many reported 99--100\% results use binary or ternary seizure detection tasks.
\item Complete five-class Bonn EEG classification is more challenging.
\item CNN-based methods can report around 92--93\% or higher.
\item This project uses a lightweight and interpretable feature-based SVM approach.
\end{itemize}

\begin{block}{Position of This Work}
The proposed method achieved 89.20\% without overlap and 90.33\% with 50\% overlap, which is competitive for a handcrafted feature-based SVM pipeline.
\end{block}
\end{frame}

\begin{frame}
\frametitle{Discussion}
\begin{columns}
\begin{column}{0.48\textwidth}
\begin{block}{Advantages}
\begin{itemize}
\item Lightweight model
\item Interpretable features
\item BO-supported parameter tuning
\item Hierarchical design targets N/F confusion
\item Strong subset classification results
\end{itemize}
\end{block}
\end{column}

\begin{column}{0.48\textwidth}
\begin{block}{Limitations}
\begin{itemize}
\item Main evaluation is window-level
\item Overlap may increase train-test similarity
\item File-level split can be tested in future work
\item Deep learning may further improve performance
\end{itemize}
\end{block}
\end{column}
\end{columns}
\end{frame}

\begin{frame}
\frametitle{Conclusion}
\begin{itemize}
\item A complete BO-SVM pipeline was developed for Bonn EEG classification.
\item Error analysis showed that N/F confusion was the main limitation.
\item A hierarchical SVM with N/F refinement was proposed.
\item Five-class accuracy improved to:
\begin{itemize}
\item 89.20\% using non-overlapping windows.
\item 90.33\% using 50\% overlapping windows.
\end{itemize}
\item Subset tasks achieved 98.67--100\% accuracy without overlap.
\end{itemize}
\end{frame}

\begin{frame}
\frametitle{Backup: Why SVM?}
\begin{itemize}
\item The Bonn dataset is relatively small.
\item SVM works well with small-to-medium datasets and handcrafted features.
\item RBF-SVM can model nonlinear decision boundaries.
\item Compared with deep learning, SVM is more lightweight and interpretable.
\end{itemize}
\end{frame}

\begin{frame}
\frametitle{Backup: Window-Level Evaluation}
\begin{itemize}
\item Each original EEG segment is divided into smaller windows.
\item Classification is performed at the window level.
\item With overlapping windows, neighboring windows share part of the signal.
\item Therefore, the 90.33\% result should be described as window-level accuracy.
\end{itemize}
\end{frame}

\begin{frame}\frametitle{}
\hspace{4cm}{\Huge Thank You}

\vspace{0.8cm}
\hspace{4.3cm}{\Large Questions?}
\end{frame}

\begin{frame}\frametitle{References}
 \begin{thebibliography}{}

\bibitem[Bonn]{Bonn}
University of Bonn.
\newblock Bonn EEG dataset for epileptic seizure analysis.

\bibitem[CortesVapnik95]{SVM}
Corinna Cortes and Vladimir Vapnik.
\newblock Support-vector networks.
\newblock {\em Machine Learning}, 20:273--297, 1995.

\bibitem[Mallat89]{DWT}
Stephane Mallat.
\newblock A theory for multiresolution signal decomposition: the wavelet representation.
\newblock {\em IEEE Transactions on Pattern Analysis and Machine Intelligence}, 1989.

\bibitem[Snoek12]{BO}
Jasper Snoek, Hugo Larochelle, and Ryan Adams.
\newblock Practical Bayesian optimization of machine learning algorithms.
\newblock {\em NeurIPS}, 2012.

\end{thebibliography}
\end{frame}

\end{document}
